\documentclass[12pt,a4paper]{article}
\usepackage{fontspec}
\setmainfont{DejaVu Serif}
\setsansfont{DejaVu Sans}
\usepackage[brazil]{babel}
\usepackage[margin=2.2cm]{geometry}
\usepackage{booktabs,tabularx,array,longtable}
\usepackage{xcolor}
\usepackage{parskip}
\usepackage{amsmath}
\usepackage{amssymb}
\usepackage{enumitem}
\usepackage{fancyhdr}
\setlength{\headheight}{15pt}
\pagestyle{fancy}
\fancyhf{}
\lhead{Análise da Complexidade do Sistema Gérard}
\rhead{Atualização C147}
\cfoot{\thepage}
\definecolor{vermelhosuave}{RGB}{194,104,104}
\begin{document}

\section*{Atualização da análise de complexidade - ciclos C36 a C147}

Esta atualização preserva a análise consolidada dos ciclos C69 a C122 e incorpora a extensão arquitetural C147, na qual os universos numéricos, os papéis, o contexto, as entidades, as pistas linguísticas e os esquemas de categoria passam a ser objetos explícitos do domínio. O sistema mantém os diagramas de Vergnaud, os diagramas complementares, as barras, os cartões e os eixos como projeções manipuláveis de um \textbf{estado semântico compartilhado}. Sobre essa base foram adicionados: atualização imediata de personagens após a curadoria; controle de funcionalidades em construção; relato contextual de falhas; memória institucional por imagens; estabilização do fluxo de persistência; curadoria semântica restrita à versão original; integração com o Gmail Web; feedback multissensorial para erro de posicionamento; ajuda contextual nas áreas do texto, de Vergnaud e do diagrama complementar; e inicialização sem conteúdo educativo com preservação da estrutura visual da tela; disponibilidade da análise qualitativa somente após posicionamento efetivo; tematização visual da janela de análise; e criação e remoção incrementais de quadradinhos sincronizadas com o estado compartilhado; refinamento do afastamento visual dos controles; e limite semântico para impedir cardinalidades superiores aos valores curados ou derivados; e refatoração das capacidades de unidades por interfaces, classe abstrata e despacho polimórfico; affordance 2,5D e pickup dos elementos; mapeamento explícito entre ordem visual e papéis semânticos; arraste com origem fantasma e mola amortecida; sincronização dos valores conhecidos do enunciado; preservação da interrogação original; controles quantitativos condicionados ao estado de Vergnaud, aos limites curados, ao piso zero e à relação aditiva; e controlador de drag-and-drop com massa, momento, amortecimento, limite de velocidade e soltura exata; e validação preventiva que permite sinal apenas na transformação ou no valor relativo e impede quantidades negativas nas medidas.

\section{Síntese do impacto}

A complexidade tecnológica geral permanece baixa a média, mas a complexidade funcional, interacional e de domínio permanece alta. Os ciclos C67 e C68 reduziram a complexidade acidental das sincronizações e eliminaram o truncamento visual. Entre C69 e C122, o crescimento ocorreu principalmente na coordenação de estados de interface, na persistência multilíngue e no scaffolding. As novas funcionalidades são localizadas e, em geral, constantes por evento, mas aumentam o número de condições que precisam preservar a fonte semântica única e não alterar a notação formal.

\begin{tabularx}{\textwidth}{p{3.8cm}p{2.8cm}X}
\toprule
\textbf{Dimensão} & \textbf{Avaliação} & \textbf{Impacto acumulado até C147}\\
\midrule
Complexidade de domínio & Alta & Papéis semânticos, sinais, incógnita e relações devem permanecer coerentes entre representações distintas.\\
Complexidade de interação & Alta & Arraste, edição, exclusão, criação de unidades, eixos, feedback de erro, painéis de ajuda, análise qualitativa e relato de falhas podem iniciar eventos distintos sobre o mesmo contexto.\\
Complexidade arquitetural & Média para alta & O estado compartilhado foi mantido e surgiram módulos específicos para suporte, curadoria multilíngue, feedback, ajuda contextual e controles incrementais das coleções. C107 reduz acoplamento ao expressar capacidades por contratos e polimorfismo; C117--C120 consolidam o estado compartilhado entre texto, Vergnaud e unidades; C121 mantém a animação de arraste fora do modelo semântico; C122 centraliza o invariante de não negatividade; C147 desloca essa regra para objetos numéricos e papéis semânticos, tornando as classificações por categoria propriedades derivadas.\\
Complexidade de renderização & Baixa a média & Coleções e barras continuam lineares no número de unidades; ícones, tooltips, tremor e painéis contextuais possuem custo constante por quadro, enquanto adicionar uma unidade redistribui a grade em tempo linear.\\
Complexidade de manutenção & Média para alta & A centralização e os contratos reduzem duplicações e condicionais concretos, mas internacionalização, persistência, Gmail e scaffolding exigem regressão cruzada entre fluxos.\\
\bottomrule
\end{tabularx}

\section{Modelo semântico compartilhado}

Para os três papéis principais das categorias atuais, o estado é representado por:
\[
S = (T, V, K, i, o, \nu),
\]
em que $T$ é a categoria, $V=(v_1,v_2,v_3)$ contém os valores, $K=(k_1,k_2,k_3)$ indica quais valores foram efetivamente modelados, $i$ identifica o papel alterado, $o$ registra a origem da ação e $\nu$ é a versão do snapshot.

A relação canônica utilizada no nível dos três papéis é:
\[
v_1 + v_2 = v_3.
\]

Essa forma cobre, entre outros casos:
\begin{itemize}[noitemsep]
  \item Parte 1 $+$ Parte 2 $=$ Todo;
  \item Estado inicial $+$ Transformação $=$ Estado final;
  \item Referido $+$ Valor relativo $=$ Referendo;
  \item Relação inicial $+$ Transformação $=$ Relação final;
  \item Transformação 1 $+$ Transformação 2 $=$ Transformação resultante.
\end{itemize}

O papel alterado define a direção da resolução. Se $v_3$ é alterado e $v_1$ é conhecido, então $v_2=v_3-v_1$. Se $v_1$ ou $v_2$ é alterado e ambos são conhecidos, então $v_3=v_1+v_2$. Quando não há papel explicitamente alterado, somente um termo ausente é derivado.

\section{Fluxo reativo}

O fluxo consolidado é:

\begin{center}
Ação do usuário $\rightarrow$ captura dos valores $\rightarrow$ estado compartilhado $\rightarrow$ resolução $A+B=C$ $\rightarrow$ snapshot $\rightarrow$ atualização de todas as representações.
\end{center}

As origens previstas são Vergnaud, diagrama complementar, eixo x, eixo vertical, edição textual, arraste, exclusão e protocolo de scaffolding. Uma trava booleana de aplicação impede que o redesenho provocado pelo snapshot seja interpretado como uma nova ação. Sem essa trava, uma atualização Vergnaud $\rightarrow$ Venn poderia acionar novamente Venn $\rightarrow$ Vergnaud, produzindo recursão ou oscilações.

\section{Complexidade temporal}

Considere $p=3$ como o número de papéis principais, $r$ como o número de representações coordenadas e $q$ como o total de unidades gráficas manipuláveis.

\subsection{Captura e resolução}

A leitura dos três elementos de Vergnaud e a resolução da relação são constantes:
\[
O(p)=O(1).
\]

No diagrama complementar, a contagem dos quadradinhos percorre as unidades. Como há número constante de regiões, o custo é:
\[
O(q).
\]

A criação do snapshot e a escolha do termo dependente permanecem $O(1)$.

\subsection{Aplicação nas representações}

Atualizar os três valores textuais de Vergnaud custa $O(1)$. Reconstruir as coleções ou barras custa $O(q')$, em que $q'$ é o número de unidades do novo estado. No ciclo C68, o cálculo do número de linhas, colunas, tamanho e espaçamento da grade é constante; a criação das unidades continua linear em $q'$. Atualizar os eixos e cartões é $O(1)$.

Na comparação de medidas, a correspondência cromática ainda exige ordenar as unidades de cada barra de baixo para cima. Para $q=r_f+r_e$:
\[
O(r_f\log r_f+r_e\log r_e)=O(q\log q).
\]

Logo, o custo total de uma transação é:
\[
O(q) \quad \text{nas coleções simples},
\]
e
\[
O(q\log q) \quad \text{na comparação com pareamento cromático}.
\]

Como as quantidades educacionais são pequenas e limitadas pela área de desenho, esse custo não representa gargalo prático.

\subsection{Custos dos ciclos C69 a C122}
Seja $m$ o tamanho textual do relato de falha, $a$ o número de áreas de ajuda e $u$ o número de opções por painel. No projeto atual, $a=3$ e $u=3$, portanto ambos são constantes.

\begin{itemize}[noitemsep]
  \item atualização de rótulos, bloqueio de campos, abertura de menu, tremor, som e tooltip: $O(1)$ por evento;
  \item geração do painel de ajuda: $O(a\cdot u)=O(1)$;
  \item serialização do relato de bug e composição da URL do Gmail: $O(m)$;
  \item cópia dos dados semânticos da versão original para uma tradução: $O(f)$, em que $f$ é o número fixo de campos curados, logo $O(1)$ no modelo atual;
  \item redimensionamento de fotografias: $O(P)$ em relação ao número de pixels processados no carregamento, sem participação nas transações matemáticas.
  \item verificação de existência de posicionamento para habilitar a análise: $O(n)$ no pior caso, em que $n$ é o número de itens arrastáveis atuais, com interrupção no primeiro item válido;
  \item aplicação do tema na janela de análise: $O(c)$ na criação, em que $c$ é o número fixo de componentes do formulário, portanto constante no modelo atual;
  \item adição de um quadradinho e redistribuição da coleção: $O(q)$, seguida da publicação do snapshot e da atualização das representações;
  \item seleção polimórfica da capacidade de adicionar ou remover: $O(1)$ por controle; a operação concreta mantém o custo $O(q)$ já determinado pela contagem e redistribuição das unidades.
  \item cálculo de habilitação de cada controle quantitativo: $O(1)$ para os três papéis e limites curados; quando a operação é executada, a reconstrução das unidades permanece $O(q)$;
  \item sincronização dos marcadores conhecidos do enunciado: $O(t)$ no número de elementos textuais manipuláveis, pequeno e limitado pelo enunciado; a interrogação original é ignorada em tempo constante por elemento.
  \item validação do sinal da transformação ou do valor relativo: $O(1)$, pois consulta no máximo as duas quantidades adjacentes, calcula um termo dependente e bloqueia a propagação quando o resultado seria negativo.
\end{itemize}

Assim, o limite assintótico das transações matemáticas continua sendo $O(q)$ nas coleções e $O(q\log q)$ na comparação. Os recursos de suporte acrescentam custo constante ou linear apenas no comprimento da mensagem.

\subsection{Custo do arraste massa--mola--amortecedor}
Durante o arraste, a posição visual $x_t$ e a velocidade $v_t$ são atualizadas em cada quadro por uma aproximação discreta:
\[
a_t = \frac{k(x^{*}_t-x_t)}{m},
\qquad
v_{t+1}=\gamma(v_t+a_t),
\qquad
x_{t+1}=x_t+v_{t+1},
\]
em que $x^{*}_t$ é a coordenada do cursor, $k$ é a rigidez, $m$ é a massa e $\gamma$ é o amortecimento. O mesmo cálculo é aplicado de forma independente aos eixos horizontal e vertical. Distância visual e velocidade são limitadas por constantes.

Cada quadro executa número constante de operações aritméticas e uma chamada de atualização visual, portanto custa $O(1)$. Para $f$ quadros de um gesto, o custo total da animação é $O(f)$. Como o temporizador usa intervalo aproximado de 16 ms, $f$ é proporcional à duração do gesto, não à quantidade de quadradinhos nem ao número de papéis semânticos. A soltura encerra o temporizador, zera a velocidade e aplica uma única atualização exata, também $O(1)$, antes da validação semântica.

O controlador não publica snapshots e não consulta a curadoria. Ele produz apenas coordenadas intermediárias para o mesmo método de movimento já utilizado pelos elementos. Assim, não altera os limites $O(q)$ e $O(q\log q)$ das transações matemáticas.



\section{Impacto específico do ciclo C68}

Antes do ciclo C68, a função de renderização limitava a quantidade visual das coleções a 12 quadradinhos. Esse limite produzia uma inconsistência observável: o estado semântico e o diagrama de Vergnaud podiam registrar, por exemplo, $8+6=14$, enquanto a coleção resultante exibia somente 12 unidades.

A correção substituiu a supressão de unidades por uma grade adaptativa. Para uma região de largura útil $L$, altura útil $H$ e quantidade $q$, o número de colunas é estimado por
\[
c=\left\lceil\sqrt{q\cdot\frac{L}{H}}\right\rceil,
\]
e o número de linhas por
\[
l=\left\lceil\frac{q}{c}\right\rceil.
\]
O tamanho e o espaçamento são então ajustados às dimensões disponíveis. O cálculo geométrico é $O(1)$ e a materialização dos $q$ quadradinhos é $O(q)$. A mudança preserva o invariante de que a quantidade de unidades gráficas deve ser exatamente igual ao valor semântico representado.


\section{Impacto acumulado dos ciclos C69 a C107}

\subsection{C69 a C72: distribuição, atualização e controle de navegação}
O ciclo C69 alterou apenas a geometria da composição de coleções: a seta foi reduzida e a resultante reposicionada. O custo de cálculo permanece $O(1)$, enquanto a renderização das unidades continua $O(q)$. O ciclo C70 reaplica nomes curados em elementos já renderizados, percorrendo uma quantidade pequena e fixa de papéis. C71 remove informação textual redundante e C72 adiciona estados desabilitados a menus incompletos. Esses ciclos aumentam pouco a complexidade computacional, mas reduzem ambiguidades e caminhos inválidos.

\subsection{C73 a C75: relato de falhas com contexto}
O formulário de bug reúne descrição, situação, categoria, idiomas e representações exibidas. O registro local e o corpo do e-mail são produzidos em tempo linear no tamanho do texto, $O(m)$. A complexidade relevante é de rastreabilidade: o relato precisa capturar o contexto sem modificar a tentativa em andamento e deve permanecer útil mesmo quando o cliente de e-mail não é aberto.

\subsection{C76 a C84: imagens e convenções de submenu}
As fotografias são carregadas e redimensionadas no momento de apresentação do submenu. O custo depende do número de pixels, mas não afeta o modelo semântico. A substituição do glifo por ícone vetorial reduz dependência de fontes. Alinhamento, largura e posição da seta são estados puramente visuais com custo constante.

\subsection{C85 a C89: persistência e curadoria multilíngue}
A correção do fechamento elimina duas gravações sucessivas e estabelece um fluxo único para o botão e o fechamento da janela. A curadoria semântica somente na versão original transforma a relação entre versões linguísticas em uma dependência explícita: traduções editam enunciado e validação, enquanto categoria, papéis, valores, sinais, participantes e termo desconhecido são herdados. O bloqueio dinâmico precisa reagir à troca do idioma, e a camada de persistência reaplica os campos originais para impedir divergências introduzidas por estados antigos da interface. A abertura do Gmail Web acrescenta codificação de parâmetros de URL em $O(m)$.

\subsection{C90: feedback multissensorial de erro}
A avaliação compara o papel semântico do número com o papel do elemento-alvo. Como o conjunto de papéis por categoria é pequeno e fixo, a decisão é $O(1)$. O tremor utiliza uma sequência curta de deslocamentos e um temporizador; o som é emitido uma vez; e o tooltip permanece associado ao elemento até que a compatibilidade seja restabelecida. O principal risco é produzir loops de animação ou deixar o aviso persistente após a correção, controlado por estado explícito de feedback.

\subsection{C91: ajuda contextual por área}
Há três áreas de ajuda e três intenções por área. A matriz de orientações possui tamanho constante e é internacionalizada em português, inglês e francês. A abertura por passagem do mouse ou clique exige controle de visibilidade, foco e fechamento ao clicar fora. Os ícones ficam no cabeçalho dos cartões, fora das formas de Vergnaud e das estruturas complementares, preservando a fidelidade representacional.


\subsection{C93--C94: inicialização sem conteúdo e preservação dos painéis}
A escolha da categoria é representada por um estado booleano ou enumerado de tamanho constante. Enquanto nenhuma categoria foi selecionada, os três contêineres visuais continuam sendo desenhados, mas seus componentes educativos internos não são materializados. A decisão de renderizar conteúdo é $O(1)$; após a seleção, os custos retornam aos limites já descritos para texto, Vergnaud e representação complementar.

A implementação intermediária C93 aplicou a condição de ocultação no nível dos painéis e produziu uma regressão visual, embora a compilação e o carregamento posterior permanecessem funcionais. C94 deslocou a condição para os filhos educativos dos contêineres. O custo computacional não se altera de modo relevante, mas a regra reduz complexidade cognitiva e de manutenção ao estabelecer uma fronteira explícita entre \textbf{estrutura persistente} e \textbf{conteúdo transitório}. A verificação passa a exigir três estados: inicial sem categoria, carregado após a seleção e retorno ou troca de contexto sem perda das divisões da tela.

\subsection{C96: análise somente após posicionamento}
A condição de disponibilidade percorre os itens arrastáveis até encontrar um elemento atualmente sobre uma área semântica de Vergnaud. O pior caso é $O(n)$, mas $n$ é pequeno e a busca termina no primeiro posicionamento válido. O botão e a rotina de abertura repetem a verificação para impedir acesso por estado indireto. A restauração recalcula a condição e oculta novamente a análise quando nenhum elemento permanece modelado. O ganho principal é metodológico: fotografias e justificativas deixam de ser produzidas antes de existir uma ação de referência.

\subsection{C97: tema visual da análise}
A aplicação da paleta do Gérard percorre uma quantidade fixa de painéis, títulos, campos e botões durante a construção da janela. O custo é constante para a estrutura atual e não participa das transações matemáticas. A separação entre tema e persistência evita que uma alteração visual modifique limites de caracteres, respostas por tentativa, protocolos automáticos ou fotografia do estado.

\subsection{C98: criação incremental de unidades}
Cada agrupamento discreto recebe um controle vetorial de adição. O clique acrescenta uma unidade, preserva os textos dos quadradinhos existentes, recalcula a grade adaptativa e publica a cardinalidade atualizada. A inserção é $O(1)$, mas a redistribuição e o redesenho percorrem a coleção, resultando em $O(q)$. Como o estado compartilhado recebe a quantidade e não uma cópia das figuras, a topologia arquitetural continua $O(r)$ em relação às representações. O controle é omitido em papéis que representam relações, transformações ou valores relativos sem coleção de unidades.

\subsection{C100: remoção incremental de unidades}
O controle vetorial de subtração é desenhado somente quando a coleção possui ao menos um quadradinho visível. A detecção do controle e a decisão de exibição percorrem as unidades do agrupamento, com custo $O(q)$ no pior caso. O acionamento remove uma unidade, preserva os textos remanescentes, recalcula a grade adaptativa e publica a cardinalidade atualizada no estado compartilhado, também em $O(q)$. A condição de visibilidade impede ações sem efeito em coleções vazias e mantém relações, transformações e valores relativos sem controles de unidades.

\subsection{C103--C104: separação visual entre controles e cardinalidades}
Os ciclos C103 e C104 alteram somente constantes geométricas usadas para posicionar os controles de adição e remoção. O cálculo da posição permanece $O(1)$ por agrupamento. O ganho é perceptivo: os botões deixam de competir com os números de cardinalidade, mantendo associação espacial com a coleção sem encobrir ou comprimir o numeral.

\subsection{C105: limite semântico das cardinalidades manipuláveis}
Antes de adicionar uma unidade, o sistema obtém o papel associado ao agrupamento, consulta o valor curado e, quando necessário, resolve o termo desconhecido pela relação $A+B=C$. Como o número de papéis principais é constante, a resolução do limite é $O(1)$. A contagem das unidades atuais continua $O(q)$; portanto, a decisão completa de permitir ou bloquear a adição permanece $O(q)$. Quando o limite é atingido, a cardinalidade não muda e o feedback multissensorial possui custo temporal constante, além do redesenho usual. A regra reduz complexidade semântica acidental ao impedir que o estado manipulável ultrapasse a fonte curada.

\subsection{C107: capacidades por contratos, herança e polimorfismo}
A refatoração separa três capacidades: consultar unidades, adicionar unidades e remover unidades. Os contratos especializados herdam do contrato-base; uma classe abstrata mantém agrupamento, papel semântico e contagem; e a implementação editável fornece as operações concretas. A verificação de capacidade e o despacho do método são $O(1)$. A contagem, a alteração e a redistribuição continuam $O(q)$, portanto o ciclo não modifica a complexidade assintótica das transações.

O ganho é de manutenção e extensão. Antes, controles e tela conheciam diretamente a classe concreta e combinavam verificações de categoria, visibilidade e cardinalidade. Depois de C107, o controle depende apenas da capacidade declarada. Uma nova representação pode ser somente consultável, apenas adicionável, apenas removível ou possuir ambas as capacidades sem exigir condicionais adicionais na tela. A herança permanece restrita ao estado realmente comum; quando não houver relação conceitual entre abstrações, a composição continua sendo a alternativa preferida.

\section{Complexidade espacial}

O estado compartilhado armazena apenas três valores, três indicadores e metadados constantes, portanto utiliza:
\[
O(1).
\]

As unidades gráficas continuam dominando a memória:
\[
O(q).
\]

Mensagens de bug e orientações contextuais utilizam espaço $O(m)$ apenas enquanto estão em memória; os estados de feedback, idioma, painel, categoria selecionada, habilitação da análise e foco dos controles de adição e remoção são constantes. Os novos quadradinhos ampliam somente o termo já existente $O(q)$. Os objetos de contrato, adaptadores e referências polimórficas são constantes por agrupamento e não alteram a ordem espacial.

A centralização não duplica os quadradinhos no modelo de domínio. O snapshot contém quantidades, não cópias das figuras.

\section{Invariantes de consistência}

Após qualquer transação, devem ser verdadeiras as seguintes condições:
\begin{enumerate}[noitemsep]
  \item todas as representações exibem o mesmo valor para cada papel conhecido;
  \item a relação $v_1+v_2=v_3$ é preservada quando há informação suficiente;
  \item o sinal do segundo termo é preservado nos números relativos;
  \item a incógnita permanece desconhecida até ser modelada ou derivada pela relação;
  \item eixos, cartões e pontos de controle refletem o mesmo valor da relação;
  \item quadradinhos e rótulos numéricos representam a mesma quantidade, sem truncamento por limite visual fixo;
  \item uma atualização de renderização não cria uma segunda ação do usuário;
  \item traduções não redefinem a semântica da situação original;
  \item a mensagem de erro desaparece quando o posicionamento se torna compatível;
  \item os ícones e painéis de ajuda não alteram as formas documentadas dos diagramas;
  \item a estrutura visual dos três painéis permanece presente mesmo quando o conteúdo educativo ainda não foi carregado;
  \item a análise qualitativa somente fica disponível quando existe ao menos um elemento atualmente posicionado no diagrama de Vergnaud;
  \item cada acionamento do controle de adição aumenta exatamente uma coleção e propaga a mesma cardinalidade ao estado compartilhado e ao diagrama de Vergnaud;
  \item o controle de remoção somente aparece em coleções com quadradinhos visíveis, reduz exatamente uma unidade e propaga a mesma cardinalidade às representações;
  \item a adição é bloqueada quando a coleção atinge o valor semântico curado ou derivado, com feedback que não modifica o snapshot;
  \item controles de unidades dependem dos contratos de capacidade, e não da classe concreta ou de condicionais por categoria;
  \item a tematização visual da janela de análise não altera respostas, protocolos, limites ou persistência;
  \item relatos de falha capturam o contexto sem modificar o estado matemático;
  \item com Vergnaud semanticamente vazio, os controles de adição e remoção permanecem inabilitados;
  \item cada adição respeita o limite curado e cada remoção respeita o piso zero, sem quebrar a relação aditiva;
  \item a curadoria é somente leitura durante a tentativa e não é sobrescrita por texto, Vergnaud, coleções ou barras;
  \item a posição visual elástica pode divergir temporariamente do cursor, mas a posição de soltura usada na validação é exatamente a coordenada final do mouse;
  \item o temporizador e o momento são encerrados no cancelamento ou na soltura e não produzem eventos sem uma interação ativa.
\end{enumerate}

\section{Impacto sobre a manutenção}

Antes do C67, correções eram adicionadas em pares: eixo vertical para barras, barras para Vergnaud, cartão para eixo, Venn para valores e valores para Venn. Esse desenho aumenta o número potencial de conexões entre $r$ representações para até:
\[
O(r^2).
\]

Com o estado compartilhado, cada representação se conecta ao coordenador, reduzindo a topologia conceitual para:
\[
O(r).
\]

O ganho é de manutenção, não necessariamente de tempo de execução. Novas representações precisam implementar duas operações: publicar alterações no estado e renderizar um snapshot. Elas não precisam conhecer todas as demais visualizações. Os ciclos posteriores acrescentam uma segunda regra de manutenção: recursos de suporte devem consultar o estado e registrar contexto, mas não podem tornar-se novas fontes da semântica. A curadoria multilíngue e o scaffolding permanecem consumidores ou mediadores do modelo compartilhado. A inicialização segue a mesma separação: os painéis pertencem à estrutura da interface, enquanto enunciados, diagramas e controles educativos são projeções condicionadas pelo estado da atividade. A análise qualitativa consulta o estado da tentativa antes de se tornar acessível, e os controles de adição publicam alterações no coordenador em vez de atualizar outras representações diretamente. Em C107, os controles passam também a depender de interfaces de capacidade. Isso reduz o número de razões para alterar a tela principal: novas representações implementam os contratos adequados e preservam o mesmo fluxo de limite, logging e sincronização. C117--C120 estendem essa separação ao enunciado e aos controles quantitativos. C121 mantém o controlador físico atrás de um contrato de arraste: qualquer elemento manipulável recebe posições visuais sem conhecer massa, rigidez ou amortecimento, e o modelo semântico continua desconhecendo a animação. C122 acrescenta uma validação anterior à publicação do snapshot: o sinal continua pertencendo à relação, enquanto os índices que representam medidas rejeitam valores negativos.


\section{Complexidade do modelo semântico explícito em C147}

A atualização C147 aumenta o número de classes, mas reduz a quantidade de decisões duplicadas nas representações. O objetivo não é modelar cada palavra ou objeto por uma subclasse. Nomes, objetos e expressões são dados de classes conceituais estáveis, como \texttt{Personagem}, \texttt{ReferenteContextual} e \texttt{PistaLinguistica}. A complexidade estrutural desloca-se de condicionais distribuídas para catálogos e objetos de valor verificáveis.

\subsection{Universos numéricos}

A criação e a validação de \texttt{NumeroNatural}, \texttt{NumeroInteiro} e \texttt{ValorDesconhecido} possuem custo $O(1)$. O papel quantitativo mantém uma referência ao universo esperado e valida a compatibilidade também em $O(1)$. O estado compartilhado continua contendo apenas três posições canônicas; portanto, a conversão de entradas antigas para objetos numéricos, a resolução de $A+B=C$ e a produção do snapshot permanecem $O(1)$ em tempo e espaço.

O ganho principal é a remoção de testes repetidos por representação. Se $r$ visualizações implementassem localmente a restrição de sinal, haveria até $r$ pontos de manutenção para cada regra. Com o objeto numérico como fonte, a regra é implementada uma vez e consumida por $r$ projeções, preservando a topologia conceitual $O(r)$ já obtida pelo estado compartilhado.

\subsection{Contexto, entidades e pistas}

Para um enunciado de comprimento $n$ e um léxico com $p$ pistas, a implementação inicial de localização por busca sequencial possui pior caso $O(pn)$. Esse custo é aceitável no corpus atual, pois $p$ e $n$ são pequenos. Se o léxico crescer significativamente, pode-se substituir a busca por trie ou autômato de múltiplos padrões, reduzindo a varredura para $O(n+z)$, em que $z$ é o número de ocorrências encontradas.

Personagens, referentes e unidades ocupam espaço linear no número de entidades $e$ e objetos contextuais $c$: $O(e+c)$. Como nomes e substantivos são dados, e não subclasses, o número de tipos Java permanece independente do tamanho do vocabulário.

\subsection{Categorias simples e compostas}

Uma categoria simples contém quantidade constante de papéis e restrições. A consulta do domínio de um papel e do papel canônico de uma posição compartilhada é $O(1)$. Para uma categoria composta com $k$ componentes e $m$ papéis, a obtenção da lista achatada custa $O(k+m)$. No caso atual, $k$ e $m$ são pequenos e limitados pelos diagramas aditivos implementados.

O padrão Composite aumenta a expressividade sem fazer da categoria a fonte do sinal. Cada nó agrega papéis que já declaram seus universos. Consequentemente, classificações legadas como ``com relação assinada'' são calculadas em $O(m)$ durante a obtenção do perfil, mas não duplicam definições matemáticas.

\subsection{Compatibilidade incremental}

As classes antigas foram preservadas como fachadas de compatibilidade:
\begin{itemize}[noitemsep]
  \item \texttt{CatalogoPapeisSemanticosAditivos};
  \item \texttt{PoliticaValoresAditivos};
  \item \texttt{CatalogoPerfisCategoriasAditivas};
  \item \texttt{PoliticaSinalTransformacaoComplementar}.
\end{itemize}
Todas delegam ao novo modelo. Essa estratégia acrescenta chamadas constantes, mas evita uma migração simultânea de todas as representações e reduz o risco de regressão funcional.

\section{Riscos e controles}

Os principais riscos introduzidos pela centralização são:
\begin{itemize}[noitemsep]
  \item derivar um termo que o desenho pedagógico deveria manter vazio;
  \item escolher o termo dependente errado quando duas regiões são alteradas na mesma ação;
  \item reconstruir componentes durante um arraste contínuo em frequência excessiva;
  \item apagar a distinção entre valor curado e valor efetivamente modelado pelo usuário;
  \item permitir que uma tradução grave valores semânticos divergentes;
  \item manter tremor, som ou tooltip após a correção;
  \item abrir painéis de ajuda acidentalmente e impedir o retorno ao diagrama;
  \item perder parâmetros do relato ao delegar o envio ao navegador;
  \item confundir ausência de conteúdo com ausência dos painéis estruturais da tela;
  \item habilitar a análise antes de existir um posicionamento ou mantê-la disponível após a restauração;
  \item criar quadradinhos duplicados, perder textos editados ou atualizar apenas uma das representações;
  \item declarar capacidades incompatíveis com a representação concreta ou usar herança apenas para reutilização acidental, sem relação conceitual estável;
  \item permitir atraso visual excessivo, oscilação prolongada ou velocidade alta que prejudiquem a precisão percebida;
  \item validar o destino usando a coordenada suavizada em vez da coordenada real de soltura;
  \item manter o temporizador ativo depois da liberação do mouse;
  \item acoplar a habilitação de vários controles quantitativos ao mesmo estado de botão ou alterar os valores curados durante a tentativa;
  \item propagar um valor relativo assinado que produza estado, referido, referendo, parte ou todo negativos, confundindo relação com cardinalidade;
  \item duplicar a definição do universo numérico em categorias ou representações, fazendo texto, Vergnaud, barras, eixo e tabuleiro aceitarem conjuntos diferentes;
  \item transformar cada palavra, nome ou objeto contextual em uma subclasse, produzindo explosão de tipos e acoplamento ao vocabulário.
\end{itemize}

Os controles implementados incluem indicadores de valor conhecido, registro do papel alterado, origem da ação, inicialização sem conteúdo com painéis preservados, dupla verificação para abertura da análise, restauração da visibilidade, preservação de textos na redistribuição das coleções, publicação da cardinalidade no estado compartilhado, trava contra recursão, fluxo único de persistência, herança semântica das traduções, estado explícito do feedback, fechamento controlado dos painéis e testes arquiteturais dos contratos, da herança especializada e do despacho polimórfico. A validação deve continuar incluindo inspeção visual e auditiva por categoria, além dos testes do modelo. Para o arraste físico, os controles incluem atraso máximo de 26 pixels, limite de velocidade, amortecimento, conclusão exata e cancelamento do temporizador. Para as quantidades positivas, testes separados verificam habilitação por agrupamento, limite superior curado, piso zero, consistência aditiva e imutabilidade dos arquivos de dados. Em C122, a suíte também reproduz o caso $6+(-8)=-2$, confirma o bloqueio antes da propagação, preserva $9+(-6)=3$ como transformação válida e verifica que categorias de relações continuam aceitando valores assinados.

\section{Verificação dos ciclos C67 a C122}

Foram executadas:
\begin{itemize}[noitemsep]
  \item compilação completa pelo Ant/NetBeans;
  \item regressão estrutural, internacionalização, vínculos, curadoria e logs;
  \item validação documental dos 113 ciclos consistentes e das três classes analíticas;
  \item compilação e renderização dos PDFs atualizados;
  \item testes específicos de estado compartilhado, grade adaptativa, bloqueio de traduções, composição do Gmail, feedback multissensorial, ajuda contextual, preservação dos três painéis no estado inicial, análise após posicionamento, tema visual da janela qualitativa, adição e remoção sincronizadas de quadradinhos, limite semântico das cardinalidades manipuláveis, contratos de unidades com herança e polimorfismo, sincronização textual com preservação da interrogação, controles quantitativos por categoria, arraste com massa, momento, reversão e soltura exata e bloqueio de sinal relativo que produziria quantidade negativa.
\end{itemize}

A inspeção visual e auditiva exaustiva em ambiente Windows continua sendo necessária para avaliar fluidez do tremor, intensidade do som, posicionamento dos tooltips e comportamento do Gmail autenticado.


\section{Verificação da atualização C147}

Foram executados:
\begin{itemize}[noitemsep]
  \item compilação completa de 242 fontes Java;
  \item 23 verificações específicas do modelo semântico explícito;
  \item testes de política de valores, perfis derivados, estado compartilhado, tabuleiro de transformação, sinal relativo, arraste físico e conclusão da modelagem;
  \item regressão estrutural completa, incluindo internacionalização, curadoria, logs, montagem, feedback e contratos de manipulação;
  \item compilação e inspeção dos PDFs atualizados.
\end{itemize}

Os testes confirmam que naturais rejeitam negativos no próprio objeto, inteiros preservam sinal, incógnitas mantêm o universo esperado, papéis determinam domínios, categorias compostas preservam os papéis dos componentes e o estado compartilhado transporta objetos numéricos sem alterar sua API pública de compatibilidade.

\section{Classificação dos ciclos C67 a C122}

C67, C73, C75, C85, C90, C91, C98, C100, C105, C107, C112, C116, C117 e C120 representam mudanças estruturais ou de generalização; C68, C70, C86, C87, C88, C94, C96, C118, C119 e C122 preservam consistência teórica e metodológica; C97, C103, C104, C109, C110, C113--C115, C121 e os demais ciclos desse intervalo concentram refinamento visual e operacional. A classificação não afirma que ciclos visuais sejam irrelevantes: eles reduzem ruído, melhoram descoberta e podem sustentar funções pedagógicas. A separação indica apenas o alcance predominante da decisão.

\section{Conclusão}

A introdução do estado semântico compartilhado aumenta a formalização do domínio, mas reduz a complexidade acidental das sincronizações locais. Até C147, o custo assintótico permanece dominado pela quantidade de unidades gráficas; persistência, feedback e ajuda acrescentam custos constantes por evento ou lineares no tamanho das mensagens. O principal benefício é a coerência: Vergnaud, Venn, barras, cartões, eixos, traduções e recursos de scaffolding deixam de competir como fontes de verdade. Cada camada passa a consultar ou modificar o modelo segundo responsabilidades explícitas, preservando a semântica e a notação formal. C96--C107 acrescentam que instrumentos qualitativos devem depender de ações observáveis, temas visuais devem permanecer desacoplados dos dados e controles de manipulação direta precisam publicar cardinalidades no estado compartilhado, tanto na adição quanto na remoção. C105 acrescenta que publicar o estado não basta: a ação precisa ser validada contra o valor semântico curado antes de alterar a coleção. C107 mostra que a reutilização dessas regras deve ocorrer por contratos de capacidade e despacho polimórfico, reduzindo dependências concretas sem alterar a complexidade assintótica. C117--C120 consolidam texto, Vergnaud, coleções e barras como projeções do mesmo estado, com curadoria imutável, limites positivos e controles independentes. C121 acrescenta custo $O(1)$ por quadro de arraste e mantém a coordenada final exata, demonstrando que uma animação com massa e inércia pode ser incorporada sem se tornar fonte de verdade ou alterar a complexidade das transações matemáticas. C122 acrescenta apenas custo $O(1)$ por mudança de sinal e torna explícito um invariante de domínio. C147 localiza esse invariante nos objetos corretos: transformações e valores relativos usam números inteiros; estados, partes, todo, referido e referendo usam naturais em $\mathbb{N}_0$; contexto, personagens e pistas são elementos semânticos reutilizáveis; e categorias simples ou compostas agregam esses elementos sem redefinir seu comportamento.

\end{document}
